\section*{Capítulo \ref{ch:b3:hydrogeology} --- Hidrogeología}

\begin{solution}{exo:b3:hydrogeology:1}
Primer tramo filtrante: $z = 120 - 40 = \qty{80}{m}$, $h = 120 - 12.5 =
\qty{107.5}{m}$, carga de presión $107.5 - 80 = \qty{27.5}{m}$. Segundo
tramo filtrante: $z = \qty{40}{m}$, $h = \qty{105.0}{m}$, carga de presión
\qty{65.0}{m}. La carga es mayor en el tramo somero, así que el agua
fluye hacia abajo, con un gradiente $(107.5 - 105.0)/40 = \num{0.0625}$:
los piezómetros se sitúan en una zona de recarga. La presión por sí sola habría
sugerido lo contrario.
\end{solution}

\begin{solution}{exo:b3:hydrogeology:2}
$q = K\Delta h/L = 12\times0.30/0.50 = \qty{7.2}{m/d}$. El área es
$\pi\times0.05^2 = \qty{7.85e-3}{m^2}$, de modo que $Q = \qty{0.0565}{m^3/d}$, o
$56.5/24 = \qty{2.36}{L/h}$. La \omterm{def:b3:hydrogeology:seepage}{velocidad de filtración} es $7.2/0.35 =
\qty{20.6}{m/d}$.
\end{solution}

\begin{solution}{exo:b3:hydrogeology:3}
$K = 2.5/86400 = \qty{2.894e-5}{m/s}$, y $k = K\mu/(\rho g) =
\num{2.894e-5}\times\num{1.306e-3}/(999.7\times9.81) =
\qty{3.85e-12}{m^2}$. A \qty{20}{\celsius}, $K$ se multiplica por
$(998.2/\num{1.002e-3})/(999.7/\num{1.306e-3}) = \num{1.30}$: $K =
\qty{3.25}{m/d}$. El cambio se debe casi por completo a la viscosidad.
\end{solution}

\begin{solution}{exo:b3:hydrogeology:4}
$T = 2.5\times180 = \qty{450}{m^2/d}$. La expansión del agua da
$\num{9e-7}\times180 = \num{1.62e-4}$, unos \qty{80}{\percent} de $S =
\num{2e-4}$; la compresión de la caliza aporta el resto.
\end{solution}

\begin{solution}{exo:b3:hydrogeology:5}
Según la \cref{prop:b3:hydrogeology:thiem}, $T = Q\ln(r_2/r_1)/(2\pi(s_1 -
s_2)) = 1500\times\ln10/(2\pi\times1.53) = \qty{359}{m^2/d}$, y $K =
359/40 = \qty{9.0}{m/d}$.
\end{solution}

\begin{solution}{exo:b3:hydrogeology:6}
$Q/4\pi T = 1200/(4\pi\times300) = \qty{0.318}{m}$. Tras 1 día, $u =
300^2\times\num{5e-4}/(4\times300\times1) = \num{0.0375}$ y $W = -0.5772
+ 3.2834 + 0.0375 - 0.0004 = \num{2.743}$: $s = \qty{0.873}{m}$. Tras 10
días, $u = \num{0.00375}$, $W = -0.5772 + 5.5860 + 0.0038 = \num{5.013}$:
$s = \qty{1.596}{m}$. Diez veces más tiempo da menos del doble de
descenso.
\end{solution}

\begin{solution}{exo:b3:hydrogeology:7}
$u \le 0.05$ cuando $t \ge r^2S/(4T\times0.05)$. A \qty{60}{m}: $3600\times
\num{2e-4}/90 = \qty{0.008}{d}$, \qty{11.5}{min}. A \qty{500}{m}:
\qty{0.556}{d}, \qty{800}{min}, más de 13 horas. El tiempo crece como $r^2$:
en una prueba de uno o dos días, los pozos de observación usados para la
recta deben estar a unos pocos cientos de metros del pozo bombeado.
\end{solution}

\begin{solution}{exo:b3:hydrogeology:8}
A mitad de camino, cada pozo está a \qty{200}{m}: $u = 200^2\times\num{2e-4}/(4
\times450\times10) = \num{4.44e-4}$, $W = \num{7.142}$, $s = 0.382\times
7.142 = \qty{2.728}{m}$ de cada uno, \qty{5.456}{m} en total. A
\qty{400}{m}, $u = \num{1.78e-3}$, $W = \num{5.757}$, $s =
\qty{2.199}{m}$. Cada pozo hace bajar el nivel del agua del otro en
\qty{2.2}{m}: su bomba debe elevar el agua \qty{2.2}{m} más arriba, lo que
cuesta energía, y el descenso que queda antes de que la toma de la bomba quede en seco
es menor.
\end{solution}

\begin{solution}{exo:b3:hydrogeology:9}
Cuatro líneas de corriente entre la base de la presa y el lecho delimitan $N_f = 5$
canales; nueve equipotenciales dividen la carga de \qty{10}{m} en $N_d = 10$
caídas de carga. $q = KHN_f/N_d = 5\times10\times5/10 = \qty{25}{m^3/d}$ por metro
de presa, y \qty{7500}{m^3/d} para \qty{300}{m}. El factor de forma
calculado, 0.512 en lugar de 0.5, daría \qty{25.6}{m^3/d} por metro.
\end{solution}

\begin{solution}{exo:b3:hydrogeology:10}
$k = na^2/8 = 0.30\times(\num{5e-5})^2/8 = \qty{9.4e-11}{m^2}$, y $K =
k\rho g/\mu = \num{9.375e-11}\times998.2\times9.81/\num{1.002e-3} =
\qty{9.2e-4}{m/s}$, \qty{79}{m/d}, dentro del intervalo de la arena limpia, entre 0.3
y \qty{150}{m/d}. El modelo omite la tortuosidad de las trayectorias, más largas
que los tubos rectos (reduce $K$), las gargantas entre los poros,
más estrechas que los poros (reduce), los poros ciegos que no
transportan flujo (reduce), y la dispersión de los tamaños de poro: como un tubo transporta un
caudal proporcional a $a^4$, las trayectorias más anchas transportan la mayor parte del agua y
una dispersión eleva $K$ por encima de la de los tubos de tamaño medio.
\end{solution}

\begin{solution}{exo:b3:hydrogeology:11}
(a)~El logaritmo se anula cuando $2.25\,Tt = r^2S$, es decir, en $r_0 =
\sqrt{2.25\,Tt/S}$. (b)~Confinado y libre:
\[
r_0 = \sqrt{\frac{2.25\times450\times30}{\num{2e-4}}} = \qty{12.3}{km},
\qquad r_0 = \sqrt{\frac{2.25\times450\times30}{0.2}} = \qty{390}{m},
\]
$\sqrt{1000} = 31.6$ veces menor. (c)~La presión
se propaga con la difusividad $T/S$, enorme cuando $S$ es pequeño: durante
la prueba, el cono alcanza fallas, los afloramientos del acuífero y los ríos,
y el descenso se aparta de la recta, lo que indica al
hidrogeólogo que se ha alcanzado un límite.
\end{solution}

\begin{solution}{exo:b3:hydrogeology:12}
(a)~Visto desde el sistema de referencia que se mueve con el agua, el frente se dispersa con
desviación típica $\sigma_x = \sqrt{2Dt} = \sqrt{2\alpha_Lvt}$. Pasa por
$x$ en torno a $t = x/v$, cuando $\sigma_x = \sqrt{2\alpha_Lx}$; a la velocidad
$v$, una dispersión de $\sigma_x$ en distancia equivale a una de $\sigma_t = \sigma_x/v$
en tiempo. (b)~Los tiempos del \qty{16}{\percent} y del \qty{84}{\percent} se sitúan una
desviación típica a cada lado de la media: $\sigma_t = (247.2 -
152.6)/2 = \num{47.3}~\text{días}$. Entonces $\alpha_L = v^2\sigma_t^2/(2x) =
0.25\times47.3^2/200 = \qty{2.8}{m}$. (La curva se calculó con
\qty{3}{m}: la aproximación normal la subestima ligeramente.) (c)~$R =
1 + 1.9\times0.32/0.30 = \num{3.03}$: el soluto viaja a un tercio de la
velocidad del agua.
\end{solution}

\begin{solution}{pb:b3:hydrogeology:1}
\textbf{1.} Lutita $363.2 - 282.5 = \qty{80.7}{m}$; caliza $282.5 -
102.5 = \qty{180}{m}$.

\textbf{2.} $h = 363.2 - 45.0 = \qty{318.2}{m}$; carga de presión en la parte superior
$318.2 - 282.5 = \qty{35.7}{m}$. El agua asciende \qty{35.7}{m} por encima de la
parte superior de la caliza, retenida por la lutita: el acuífero es confinado.

\textbf{3.} A partir de $h_2 - h_1 = 550a - 100b = -1.62$ y $h_3 - h_1 = 250a
- 600b = -1.44$: $a = \num{-2.715e-3}$ y $b = \num{1.269e-3}$ por
metro.

\textbf{4.} $|\nabla h| = \sqrt{a^2 + b^2} = \num{0.0030}$. El agua
fluye según $-\nabla h = (\num{2.715e-3}, \num{-1.269e-3})$, hacia el este y
el sur: acimut $\ang{180} - \arctan(2.715/1.269) = \ang{115}$.

\textbf{5.} Los puntos anteriores a unos \qty{10}{min} quedan por encima de la recta: allí
$u > 0.05$ y la curva de Theis está por encima de su aproximación logarítmica
($W > -\gamma - \ln u$, como muestra la figura de la función de pozo).

\textbf{6.} La mejor recta que ajusta los puntos desde \qty{15}{min} en adelante pasa por
\qty{1.75}{m} a \qty{100}{min} y \qty{2.63}{m} a \qty{1000}{min}:
$\Delta s = \qty{0.87}{m}$ por ciclo logarítmico.

\textbf{7.} $T = 2.3\times2160/(4\pi\times0.872) = \qty{453}{m^2/d}$.

\textbf{8.} La recta desciende \qty{0.87}{m} por ciclo desde \qty{1.75}{m} a
\qty{100}{min}: el cero se alcanza $1.75/0.872 = 2.01$ ciclos antes,
$t_0 = \qty{0.98}{min} = \qty{6.83e-4}{d}$.

\textbf{9.} $S = 2.25\,Tt_0/r^2 = 2.25\times453\times\num{6.83e-4}/3600 =
\num{1.9e-4}$, dentro del intervalo de los \omterm{def:b3:hydrogeology:aquifer-types}{acuíferos confinados}.

\textbf{10.} A \qty{15}{min}, $u = 3600\times\num{1.9e-4}/(4\times453
\times0.0104) = \num{0.036}$, y menor después.

\textbf{11.} $K = 453/180 = \qty{2.5}{m/d}$, una caliza fracturada.

\textbf{12.} $u = 3600\times\num{2e-4}/(4\times450\times2) =
\num{2.0e-4}$, $W = -0.5772 + 8.5172 + 0.0002 = \num{7.940}$, $s =
0.382\times7.940 = \qty{3.03}{m}$, como se registró.

\textbf{13.} $u = 10^6\times\num{2e-4}/(4\times450\times30) =
\num{0.0037}$, $W = \num{5.02}$, $s = \qty{1.92}{m}$.

\textbf{14.} $r_0 = \sqrt{2.25\times450\times30/(\num{2e-4})} =
\qty{12.3}{km}$, más que la anchura de la región. El cono ha alcanzado
los bordes de la caliza, sus afloramientos y la falla, que quizá lo
enfrenten a roca menos permeable: el acuífero infinito de Theis ya no
lo describe, y solo el registro de los descensos dirá si un
límite alimenta el pozo o detiene el flujo.

\textbf{15.} $2160\times30 = \qty{64800}{m^3}$. Procede del almacenamiento:
el agua liberada por la expansión del agua y la compresión de
la caliza, en toda la superficie del cono. Integrando $S\,s$ sobre
el cono recupera $Qt$.

\textbf{16.} $u = 10^6\times0.18/(4\times450\times30) = \num{3.33}$, $W
= \num{0.0086}$, $s = \qty{3}{mm}$: la granja no lo notaría. El
bombeo tomaría agua de los poros cercanos al pozo en lugar de la presión
de todo el acuífero.

\textbf{17.} $q = 2.5\times0.0030 = \qty{0.0075}{m/d}$; $v = 0.0075/0.18
= \qty{0.042}{m/d}$, \qty{15}{m} al año.

\textbf{18.} $200/0.0417 = \num{4800}\text{ días}$, \num{13.1}~\text{años}; el
contaminante, \num{32.9}~\text{años}.

\textbf{19.} A través de un cilindro de radio $r$ el caudal $Q$ atraviesa el
área $2\pi rb$ de la que los poros son la fracción $n_e$: $\mathrm dr/\mathrm
dt = -Q/(2\pi rbn_e)$, e integrando $2\pi bn_er\,\mathrm dr = -Q\,
\mathrm dt$ desde $r$ hasta 0 da $t = \pi n_ebr^2/Q$. Aquí $\pi\times0.18
\times180\times200^2/2160 = \num{1885}~\text{días}$, \num{5.2}\text{ años} para el
agua y $2.5\times5.2 = \num{12.9}\text{ años}$ para el contaminante: el pozo
lo hace llegar más de dos veces antes.

\textbf{20.} $\alpha_L = \qty{20}{m}$ y $\sqrt{2\times20\times200} =
\qty{89}{m}$. La concentración alcanza el \qty{16}{\percent} de su
valor final a \qty{89}{m} por delante del frente medio: ese nivel llega al
pozo cuando el frente medio aún está a \qty{89}{m} de distancia, en $(111/200)^2\times12.9 = \num{4.0}\text{ años}$, un tercio
del tiempo medio. Las primeras trazas importan para el agua potable, no la
media.

\textbf{21.} A favor: el pozo capta el contaminante, así que no puede extenderse
más aguas abajo, y el agua puede tratarse donde se bombea.
En contra: el bombeo arrastra el contaminante hacia un pozo de abastecimiento antes de lo que
recorrería de otro modo, y el agua debe vigilarse y
tratarse durante décadas.

\textbf{22.} La \omterm{def:b3:hydrogeology:storage}{transmisividad} de la caliza Varden es $T \approx
\textbf{\qty[text-series-to-math=true]{450}{m^2/d}}$ (\qty{453}{m^2/d}
a partir del ensayo).
\end{solution}
